\documentclass[11pt]{article}
\usepackage{amsmath, amssymb, amsthm}
\usepackage{geometry}
\geometry{margin=1.2in}
\usepackage{hyperref}

\newcommand{\R}{\mathbb{R}}
\newcommand{\E}{\mathbb{E}}
\newcommand{\norm}[1]{\left\|#1\right\|}
\newcommand{\abs}[1]{\left|#1\right|}

\title{Implicit Bias of GD vs.\ Adam for Diagonal Linear Networks}
\date{}

\begin{document}
\maketitle

\section{Setup}

Consider a two-layer diagonal linear network with parameters $\mathbf{u}, \mathbf{v} \in \R^d$, whose
prediction on input $\mathbf{x} \in \R^d$ is
\[
    f(\mathbf{x}) = \mathbf{x}^\top \boldsymbol{\beta}, \qquad \boldsymbol{\beta}_i = u_i v_i.
\]
We train on $n$ samples $\{(\mathbf{x}_j, y_j)\}_{j=1}^n$ with the squared loss
\[
    \mathcal{L} = \frac{1}{n} \sum_{j=1}^n \bigl(\mathbf{x}_j^\top \boldsymbol{\beta} - y_j\bigr)^2,
\]
initialised symmetrically at small scale: $u_i = v_i = c$ for small $c > 0$.

The ground-truth teacher is \emph{sparse}: $\boldsymbol{\beta}^* \in \R^d$ has support of size $\rho d \ll d$.

\section{Implicit Bias of Gradient Descent}

The gradient-flow equations for $\mathbf{u}$ and $\mathbf{v}$ are
\[
    \dot{u}_i = -\frac{\partial \mathcal{L}}{\partial \beta_i} \, v_i, \qquad
    \dot{v}_i = -\frac{\partial \mathcal{L}}{\partial \beta_i} \, u_i.
\]
With symmetric initialisation $u_i(0) = v_i(0) = c$, the conservation law
$\tfrac{d}{dt}(u_i^2 - v_i^2) = 0$ ensures $u_i = v_i$ throughout training, so
$\beta_i = u_i^2 \geq 0$.  Writing $g_i = \partial \mathcal{L}/\partial \beta_i$, the
effective dynamics on $\boldsymbol{\beta}$ become
\[
    \dot{\beta}_i = -2g_i \,\beta_i^{1/2} \cdot \beta_i^{1/2} = -2g_i\,\beta_i.
\]
Equivalently, the gradient flow on the \emph{square-root} parametrisation
$\phi_i = \sqrt{\beta_i} = u_i$ is
\[
    \dot{\phi}_i = -g_i \,\phi_i.
\]
This is precisely \textbf{mirror descent} with the Bregman potential
$\Phi(\boldsymbol{\phi}) = \sum_i \phi_i^2 \log \phi_i$, which in the original
$\boldsymbol{\beta}$ coordinates corresponds to an $\ell_{1/2}$-like regulariser.  As
$c \to 0$, gradient flow is therefore biased towards the \textbf{minimum
$\ell_1$-norm interpolant}~\cite{woodworth2020kernel}, which recovers the
sparse teacher whenever $n$ exceeds a threshold of order $\rho d \log(d/\rho d)$.

\paragraph{Intuition.}  The factor of $v_i$ in $\dot{u}_i$ creates a
\emph{rich-get-richer} dynamic: coordinates that are already large attract
larger updates and grow faster, while small coordinates remain small.  This
self-amplification is the mechanism that selects a sparse solution.

\section{Why Adam Fails}

The Adam update for $u_i$ at step $t$ is
\[
    u_i^{(t+1)} = u_i^{(t)} - \eta \,\frac{\hat{m}_i^{(t)}}{\sqrt{\hat{v}_i^{(t)}} + \epsilon},
\]
where $\hat{m}_i \approx g_i v_i$ and $\hat{v}_i \approx \E[(g_i v_i)^2]$ are the
bias-corrected first and second moment estimates.  The effective per-coordinate
learning rate is therefore
\[
    \eta_i^{\text{eff}} \approx \frac{\eta}{\sqrt{\E[g_i^2 v_i^2]}}.
\]
Crucially, this \textbf{normalises away the $|v_i|$ factor}: regardless of how
large or small $v_i$ is, the update is rescaled to have roughly unit magnitude.
The rich-get-richer amplification is destroyed, and all coordinates receive
approximately equal effective step sizes.

Without this self-amplification, Adam converges quickly to a \textbf{dense
interpolating solution} that spreads weight across all $d$ coordinates.  Since
the true signal is sparse and $n \ll d$, this dense solution has poor test
error.

\section{Why Test Error Increases with $\alpha = n/d$ for Adam}

In the underdetermined regime ($n < d$), the loss landscape contains infinitely
many interpolating solutions.  Adam effectively selects the one closest (in a
distorted, coordinate-wise metric) to the initialisation, which is a dense,
low-magnitude vector.  As $\alpha$ increases, the additional training
constraints force this dense solution to satisfy more equations, pushing it
further from the sparse ground truth.  Formally, for a sparse $\boldsymbol{\beta}^*$
the minimum-$\ell_2$-norm interpolant satisfies
\[
    \boldsymbol{\beta}^{\ell_2} = \mathbf{X}^\top (\mathbf{X}\mathbf{X}^\top)^{-1} \mathbf{y},
\]
and its test error $\E_{\mathbf{x}}\bigl[(\mathbf{x}^\top(\boldsymbol{\beta}^{\ell_2} -
\boldsymbol{\beta}^*))^2\bigr]$ is non-monotone in $n$ for sparse signals: in the
regime $\rho d \ll n \ll d$, each additional constraint reshapes the dense
solution in a direction that need not reduce its distance from $\boldsymbol{\beta}^*$,
and can increase it.  This is a manifestation of the \emph{harm of
interpolation} for mismatched inductive bias~\cite{koehler2021uniform}.

\section{Summary}

\begin{center}
\begin{tabular}{lll}
\hline
Optimizer & Implicit bias & Generalisation \\
\hline
Gradient descent & Minimum $\ell_1$ norm (sparse) & Good: recovers sparse teacher \\
Adam & Dense interpolant (no sparsity) & Poor: test MSE flat or increasing in $\alpha$ \\
\hline
\end{tabular}
\end{center}

The generalisation gap is \emph{not} a tuning problem.  No choice of Adam
learning rate can restore the $\ell_1$ implicit bias, because the bias is
structural: it arises from the interaction between the gradient and the current
parameter magnitude, which Adam's normalisation removes.

\begin{thebibliography}{9}
\bibitem{woodworth2020kernel}
    B.\ Woodworth, S.\ Gunasekar, J.\ D.\ Lee, E.\ Moroshko, P.\ Savarese,
    I.\ Golan, D.\ Soudry, N.\ Srebro.
    \textit{Kernel and Rich Regimes in Overparametrized Models}.
    COLT 2020.

\bibitem{koehler2021uniform}
    F.\ Koehler, L.\ Zhou, D.\ J.\ Sutherland, N.\ Srebro.
    \textit{Uniform Convergence of Interpolators: Gaussian Width, Norm Bounds,
    and Benign Overfitting}.
    NeurIPS 2021.
\end{thebibliography}

\end{document}
